\section{Proof-directed higher-order sketch decomposition}
\label{sec:sketches}
\status{Integration design, not an implemented synthesizer. The archive includes a
hand-written, uncompiled \code{List}/\code{Int} fold specimen. It does not close
Leant's original Church-encoded indexing query.}

\subsection{The specific obstruction in Leant and Djex}
The inspected Leant report describes an unusually useful diagnostic. Under the
original bounded Haskell Exference indexing query, all 256 inspected candidates
were behaviorally falsified. Yet a suitable integer-case step synthesized in
isolation at candidate 5, and a corresponding step in a larger ambient context
appeared at candidate 219. Function-valued folds also appeared separately.
These observations locate a composition and ordering problem, not simply a
missing source-language construct~\cite{leant-indexing}.

The report explicitly warns that candidate 219 is not the internal ordinal of
that step in the original fold search. Partial terms, pending goals, usage
history, and context differ. Accordingly, the proposal here is not to infer a
magic ranking penalty from those two numbers. It is to make a proof schema
expose smaller, semantically specified synthesis problems before complete
candidate terms are enumerated.

\subsection{The proof rule should determine the holes}
For a right fold, choose a carrier type $B$ and a relation
\[
 \mathcal R:\operatorname{List}(A)\to B\to\operatorname{Prop}.
\]
The two proof obligations are
\begin{align}
 &\mathcal R([],b_0),\label{eq:fold-base}\\
 &\forall a,\ell,b,\quad
   \mathcal R(\ell,b)\Longrightarrow
   \mathcal R(a::\ell,\operatorname{step}(a,b)).\label{eq:fold-step}
\end{align}
They imply
\begin{equation}
\label{eq:fold-rule}
 \forall\ell,\quad
 \mathcal R\bigl(\ell,\operatorname{foldr}(\operatorname{step},b_0,\ell)\bigr)
\end{equation}
by structural induction. Forge can register this as a reusable proof theorem.
The synthesis task is then to construct $b_0$ and \code{step} satisfying these
contracts, not an entire recursive program from a bare outer type.

The relation is the key. Merely choosing the fold combinator and leaving its
holes constrained only by types reproduces the original ambiguity. Conversely,
choosing an arbitrary expressive relation and asking a solver to discover it is
not yet an algorithm. A first implementation should use an explicit inventory
of specification templates: equality with a reference fold, a size equation, a
pointwise relation for function carriers, a predicate preserved by constructors,
or a packet of arithmetic invariants discovered by Section~\ref{sec:ideals}.

\subsection{Indexing explains why the carrier must be a function}
Let $d:A$ be a default. Define a total reference operation on ordinary lists and
mathematical integers by
\[
 \operatorname{at}(d,[],i)=d,
\]
\[
 \operatorname{at}(d,a::\ell,i)=
 \begin{cases}
 d,&i<0,\\
 a,&i=0,\\
 \operatorname{at}(d,\ell,i-1),&i>0.
 \end{cases}
\]
The source list is the structural argument, while the index changes during
processing. Lambda-lifting the index into the fold result suggests
\[
 B=\Z\to A,
 \qquad \mathcal R(\ell,k)\equiv
       \forall i,\ k(i)=\operatorname{at}(d,\ell,i).
\]
The base hole has the contract $\forall i,b_0(i)=d$. The step hole has the
contract
\begin{equation}
\label{eq:index-contract}
 \operatorname{step}(a,k)(i)=
 \begin{cases}
 d,&i<0,\\
 a,&i=0,\\
 k(i-1),&i>0.
 \end{cases}
\end{equation}
The simple solution is therefore
\begin{lstlisting}
base = fun i => default
step = fun head continuation i =>
    if i < 0 then default
    else if i = 0 then head
    else continuation (i - 1)
answer = foldr step base xs index
\end{lstlisting}

When an approved integer-case primitive already has exactly this contract, as in
the Leant diagnostic, it should be tried as a whole component before synthesizing
its branch structure. This is a concrete use for Djex: solve a small typed
component problem under an explicit semantic contract, then return the candidate
through its existing source-typed boundary. Forge supplies the fold theorem and
the proof obligations that the type-directed search alone does not express.

The carrier heuristic is also specific: when a non-structural argument changes
in a recursive specification and the result type is $R$, try the corresponding
function type $I\to R$ as a fold carrier. For several changing arguments, try a
curried or product carrier with the original dependency order. This is an OR
alternative, not a theorem that all such functions are folds.

\subsection{Separate executable context from specification context}
The step program in \eqref{eq:index-contract} needs $d,a,k,i$, but not the
original whole list or the original outer index. The proof of its contract may
still quantify over the tail list $\ell$ and use
$\mathcal R(\ell,k)$. Thus there are two contexts:
\[
 \Gamma_{\mathrm{code}}=(d,a,k,i),
 \qquad
 \Gamma_{\mathrm{proof}}=\Gamma_{\mathrm{code}},\ell,
                         h:\mathcal R(\ell,k).
\]
A ghost variable may occur in the specification without being available as a
program component. Erasing it from the proof context would lose the induction
hypothesis; making it available as an executable value could let the synthesizer
bypass the intended fold by calling a reference function on the tail.

Context projection therefore operates on the executable inventory and closes
under all type dependencies. If a retained variable's type mentions another
local, that dependency is retained or abstracted explicitly. The full source
context remains available to the proof checker. The small-context search is a
particular grammar restriction, not a globally sound reason to prune every
larger-context search branch. A separate broader alternative can remain queued.

This proposal sharpens the diagnosed context effect without claiming that the
specific Leant miss has already been traced to one variable or one penalty.
It also makes an implementation experiment precise: retain the original outer
query, but replace the internal monolithic branch with the registered fold
schema and separately constrained holes.

\subsection{A concrete search protocol}
\begin{lstlisting}
prepare_specification(goal):
    recover source function, recursive argument, result, and laws
    retrieve approved recursor/fold schemas by input and result shape
    propose carriers by dependency-preserving lambda lifting
    for each schema and carrier:
        instantiate a registered relation template
        derive base and step contracts by applying the schema theorem
        compute executable inventories for the holes
        retain ghost variables and assumptions in proof contexts
        enqueue this plan as an alternative

solve_plan(plan):
    request typed candidates for each hole from Djex or native search
    reject concrete counterexamples where evaluation is available
    prove the candidate's whole hole contract in Lean
    retain a universally quantified hole theorem, not a sample verdict
    apply the schema theorem to assemble the original proof
    replay at the exact original goal and source inventory
\end{lstlisting}

The initial component grammar should contain local values, constructors,
projections, applications of approved providers, and the selected integer-case
primitive. For first-order arithmetic holes, the constructive projection worker
can replace candidate enumeration. More general finite arithmetic grammars can
be sent to cvc5's SyGuS interface as proposal problems, followed by Lean proof
of the returned candidate's contract. The documented cvc5 API provides
\code{synthFun}, \code{addSygusInvConstraint}, and synthesis-result extraction;
no cvc5 adapter was executed in this work~\cite{cvc5}.

The relation to refinement-type synthesis is direct: semantic types decompose
the specification before full programs are built. Synquid is a primary example
of this approach. TYGAR provides a complementary example of using an abstract
type graph and refinement to guide component composition~\cite{synquid,tygar}.
The proposed Forge specialization is to use already proved Lean recursor laws
as the decomposition authority and to make the resulting hole theorems reusable
in the original proof. It is not a new claim of decidability for dependent or
impredicative synthesis.

\subsection{Church encodings require an explicit representation law}
\label{sec:church}
The original diagnostic uses a Church-encoded list, not Lean's inductive
\code{List}. A theorem about \code{List.foldr} alone does not close that query.
To transfer the sketch, require a relation such as
\[
 \operatorname{Rep}(c,\ell)\Longrightarrow
 c\ B\ \operatorname{step}\ b_0
 =\operatorname{foldr}(\operatorname{step},b_0,\ell),
\]
at the exact carrier and universe selection in use. For an explicit encoding
of an ordinary list, this is a straightforward fold law. For an arbitrary
polymorphic value, it must not be inferred merely from its displayed type.
Parametricity, representability, totality, and available eliminations are
substantive assumptions, especially across Haskell and Lean.

The source Haskell \code{Int} is not automatically the same object as Lean's
unbounded \code{Int}. The actual branch primitive, source binder order,
polymorphic selections, and dictionary payloads must be retained. The archive's
\code{IndexingSketch.lean} intentionally proves only a proposed equation between
two ordinary-list definitions at mathematical integers. It is hand-written and
uncompiled, and it is not substituted for acceptance of the original query.

\subsection{What should be measured in the real integration}
Record the number of schema alternatives, candidate holes, full contract checks,
reusable hole theorems, and reconstructed whole programs. Compare the same outer
query, source inventory, budgets, and acceptance predicate. An isolated hole
success is diagnostic evidence, not a solved original benchmark. Preserve a
negative control that actually reaches predicate checking, rather than calling a
zero-candidate timeout a successful rejection.

The important hypothesis is that proof-directed decomposition reduces wasted
whole-program verification and context-induced branching. This hypothesis is
plausible from the inspected diagnostic but remains unmeasured here. It should
be tested before changing Djex's global search penalties or claiming a broader
Church-encoding milestone.

\section{A cross-worker proof, with a precise division of labor}
\label{sec:composition}
Consider a full binary tree $t$, with $L=L(t)$ leaves and $N=N(t)$ internal
nodes. A goal combining structural and arithmetic reasoning is
\begin{equation}
\label{eq:combined-goal}
 \exists y\in\Z,\quad
 N\le y<N+15
 \ \land\ 2y\equiv L-N+1\pmod6
 \ \land\ y\equiv N\pmod5.
\end{equation}
Here natural counts are explicitly cast to integers.

The finite-algebra worker, using counts modulo six, proves
$L-N\equiv1\pmod6$ for all trees. Hence $L-N+1\equiv2\pmod6$, supplying the
divisibility guard required to normalize the first arithmetic congruence.
The integer worker reduces it to $y\equiv1\pmod3$, merges it with
$y\equiv N\pmod5$, and chooses the first such integer at least $N$. Its period
is 15, so the chosen value is strictly below $N+15$.

For example, with a normalized residue representative
\[
 r(N)=\bigl(1+3(2(N-1)\bmod5)\bigr)\bmod15,
\]
one possible displayed witness is
\[
 w(N)=N+((r(N)-N)\bmod15).
\]
The factor 2 is the inverse of 3 modulo 5. This is a hand-simplified explanatory
witness; the generic prototype retains its own instruction graph rather than
claiming that the simplification was discovered automatically.

The structural component supplies a proved property of all trees. The arithmetic
component supplies one uniform witness and proves its constraints. Neither
worker invents assumptions to make the other succeed. The existing Forge
controller only needs to route the missing divisibility demand to the finite
summary worker and then introduce the resulting proof. This composition is a
specified end-to-end target, not a Lean run in the archive.
